\needspace{20\baselineskip}
\section{DMA buffer：配置、cache 與記憶體屬性}
\label{sec:dma}

\begin{oneline}
PL 讀寫的 DDR buffer 由 driver 透過 DMA API 配置，kernel 同時給出 CPU 用的位址與 PL 用的 DMA 位址，並讓 CPU 以 Normal non-cacheable 屬性存取它；在同一組開機檔上，這個屬性就是 driver 版本比 UIO 版本快約 1.1\% 的原因。
\end{oneline}

\begin{before}
device tree 保留一塊固定的 256 MB，標成 \code{no-map}（kernel 完全不碰）：

\snip[{}]{v100_dt_reserved}

app 再經 \code{/dev/mem}（第 148 行），用寫死的實體位址 \code{0x6000_0000}（第 140 行的 \code{FUSE_ARENA_PA_DEFAULT}）把整塊 \code{mmap} 進來：

\snip{v100_uio_arena_init}

問題：位址寫死在程式與 device tree 兩個地方；256 MB 只用到約 6 MB，其餘 kernel 也不能用；\code{/dev/mem} 能讀寫\textbf{任何}實體位址（第 10 章）；映射出來的屬性是 Device-nGnRnE，寫入不能合併、不能非對齊存取，連 \code{memset} 都不能用。
\end{before}

\begin{tech}
\begin{itemize}
\item \code{dma_set_mask_and_coherent(dev, DMA_BIT_MASK(40))}：告訴 kernel 這個裝置能定址 40 bits。
\item \code{dmam_alloc_coherent()}：配一塊 CPU 與裝置「一致」（coherent）的記憶體，回傳 CPU 位址與 \code{dma_addr_t}；\code{dmam_} 表示 driver 離開時自動歸還。
\item CMA（Contiguous Memory Allocator，連續記憶體配置器）：開機時劃出的一塊區域（這裡 256 MB），平時借給可以搬走的一般用途，driver 要大塊實體連續的記憶體時再清出來。
\item \code{dma_addr_t}：裝置看到的位址。沒有 IOMMU 時等於實體位址；有 SMMU 時是 I/O 虛擬位址（IOVA，第 11 章），程式不該假設兩者相同。
\end{itemize}
\end{tech}

\begin{gain}
\begin{itemize}
\item 位址不再寫死：每次載入由 kernel 決定，driver 再寫進硬體。driver 只配 8 MB（module 參數 \code{arena_mb}），app 需要 6 MB；256 MB 保留區拿掉了（第 10 章）。
\item 卸載時自動歸還：沒有 SMMU 時連續載入 10 次、每次 64 MB（共 640 MB，CMA 只有 256 MB），10 次都成功。
\item CPU 端屬性從 Device 變成 Normal non-cacheable，寫入可以合併：CPU 寫 buffer 快 3 到 5 倍，pipeline 的「寫輸入」從 \splitDevIn{} ms 縮到 \splitNcIn{} ms（每次執行印出的最後一幀，各 3 次），fps 高 1.2\%（開機組 D，只差這個屬性的兩組）。可以 \code{memset} 這一點是依 ARM 規格推得，app 沒有改用，也沒有在板上測。
\end{itemize}
\end{gain}

\subsection{實際例子}

\textbf{probe 的第 3 步。}先設 DMA mask，再配 buffer，最後檢查對齊：

\snip{drv_probe_dma}

\code{DMA_BIT_MASK(40)} 對應 CSR 的 40-bit 位址暫存器；kernel 配記憶體時保證位址在這個範圍內。\code{dmam_alloc_coherent} 的回傳值 \code{td->arena} 是 CPU 位址（kernel 用），\code{&td->arena_dma} 收到 DMA 位址（寫進 PL 的暫存器）。64 KB 對齊是因為 app 的配置以 64 KB 為單位（head 的兩個輸入 slot 相隔 \code{0x10000}）。

\textbf{8 MB 裡放什麼。}buffer 的配置由 app 決定，driver 只管「給一塊、檢查 offset 不越界」。\code{FUSE_ARENA_NEED} 是 6 MB，後面還有一個 \code{_Static_assert}，配置長大超過 6 MB 就編譯失敗：

\snip{app_arena_map}

依序是：backbone 權重、兩個 activation slot（17 個 block 輪流使用）、輸入影像、backbone 輸出；\code{0x400000} 起是 head 的權重、兩個輸入 slot、三層中間結果與兩個輸出 slot。

\textbf{板上：buffer 從哪裡來，取決於有沒有 SMMU。}關掉 SMMU 時（開機組 I），DMA 位址 \code{0x6a10_0000} 落在 CMA（從 \code{0x6a00_0000} 起 256 MB）裡，\code{CmaFree} 少了 8 MB 多；連續 10 次各要 64 MB，超過 CMA 的總量也都成功，表示每次卸載都真的歸還了：

\boardfile{I1_driver_state_nosmmu.log}{1,19-24,56-60}

發佈映像打開了 SMMU（開機組 H）：DMA 位址變成 IOVA \mbox{\code{0xff_ff80_0000}}，\code{CmaFree} 幾乎沒少，因為 SMMU 可以把分散的實體頁排成連續的 IOVA，DMA API 改從一般記憶體配頁（本章「原理」）：

\boardfile{H10_driver_state.log}{19-24}

\begin{pit}
一開始用 \code{/proc/meminfo} 的 \code{CmaFree} 檢查有沒有漏記憶體，卸載後數字卻沒有回升，看起來像漏掉了。其實不是：CMA 沒人用時會借給 page cache（kernel 快取檔案內容的記憶體）這類可搬移的用途，\code{CmaFree} 會隨讀檔變化，不能當證據。後來改用「載入總量超過 CMA」的方法，才得到能下結論的證據。
\end{pit}

\textbf{量測：記憶體屬性值多少。}driver v1.0 剛完成時比 UIO 版本快約 0.8\%（開機組 A，第 13 章），但當時一起改了三件事：等硬體的地方、buffer 位置、CPU 端屬性。在 v1.0.0 的開機檔上（開機組 D）做四組實驗，每組 pipeline 跑 6 次、分兩輪且順序相反：

\begin{center}
\small
\begin{tabularx}{\textwidth}{l l l l X}
\toprule
組別 & 等硬體的地方 & buffer 位置 & CPU 端屬性 & 跟上一組的差別 \\
\midrule
\code{uio} & app & \code{0x6000_0000} & Device-nGnRnE & 起點（v1.0.0） \\
\code{uio_68} & app & \code{0x6800_0000} & Device-nGnRnE & buffer 位置 \\
\code{drv_dev} & driver & CMA & Device-nGnRnE & 等硬體的地方與 buffer 來源（一起換） \\
\code{drv_nc} & driver & CMA & Normal non-cacheable & 只有 CPU 端屬性 \\
\bottomrule
\end{tabularx}
\end{center}

\code{drv_dev} 靠 driver 的 module 參數 \code{map_device=1}：\code{twn_mmap()} 改用 \code{/dev/mem} 那種 Device 屬性映射（第 4 章）。

\boardfile[，第 27--31 行：每組 6 次的平均]{D1_ab_measure.log}{1-2,27-31}

前三組都在 228.8 到 229.1 fps，只有 \code{drv_nc} 跳到 \abDrvNc{} fps。直接量 CPU 存取 buffer 的速度（\code{twn_bench} 連續寫、讀 1 MB），寫入差很多，讀取一樣：

\boardfile[：第 11 行是 Normal non-cacheable，第 20 行是 \code{map_device=1}]{D2_bench.log}{11-12,20}

pipeline 開 \code{PIPE_TIMING=1}，每次執行（120 幀）印出最後一幀每一步的時間：只有 \code{write-in}（CPU 把輸入影像寫進 buffer）變了，從 \splitDevIn{} ms 變成 \splitNcIn{} ms，和這三次執行平均的每幀週期差（約 0.04 ms）相當：

\boardfile[：前三次是 Normal non-cacheable，後三次是 \code{map_device=1}]{D6_stage_times.log}{3-21}

\subsection{原理}

\textbf{為什麼會有 cache 問題。}CPU 讀寫 DDR 會經過 cache；PL 經 HP0--HP3 port（PL 存取 DDR 的四個高速 AXI 埠）直接到 DDR 控制器，\textbf{不會}去查 CPU 的 cache。buffer 若是 cacheable：CPU 寫的輸入可能還停在 cache，PL 讀到舊資料；PL 寫進 DDR 的輸出，CPU 可能命中 cache 裡上一幀的內容。處理方法有三種：

\begin{center}
\small
\begin{tabularx}{\textwidth}{l X X}
\toprule
做法 & 怎麼做 & 代價 \\
\midrule
coherent buffer（這次用的） & \code{dma_alloc_coherent}：CPU 端映射成 non-cacheable，讀寫直接到 DDR & CPU 存取比 cacheable 慢；app 只寫輸入、讀輸出，量不大 \\
streaming DMA & 一般 cacheable 記憶體；PL 讀之前 \code{dma_sync_single_for_device}（寫回 cache），PL 寫完 \code{dma_sync_single_for_cpu}（作廢 cache） & 每幀都要做 cache 維護，app 也要配合 \\
硬體一致性 & 改用 HPC port，PL 的存取經過 CCI（會查 CPU cache 的互連）；AXI master 送 cacheable、shareable 的屬性；節點加 \code{dma-coherent} & 要改 RTL 與 block design、重建 bitstream \\
\bottomrule
\end{tabularx}
\end{center}

kernel 不知道 HP port 是什麼，它看 device tree：節點沒有 \code{dma-coherent}，就把裝置當成 non-coherent，\code{dma_alloc_coherent} 給出 non-cacheable 映射。

\textbf{\code{dma_alloc_coherent} 實際做了什麼。}沒有 SMMU 時：從 CMA 拿實體連續的頁（一般的配置器一次最多給 4 MB，而且久了記憶體會零碎，不保證拿得到）；把這些頁在 cache 裡的內容寫回並作廢；另外映射一個 Normal non-cacheable 的 kernel 位址回傳；DMA 位址等於實體位址。有 SMMU 時：PL 看到的是 IOVA，實體頁不必連續，DMA API 從一般記憶體配頁，再由 SMMU 排成一段連續的 IOVA，所以不佔 CMA。driver 的程式兩種情況完全一樣。

\textbf{為什麼不用 \code{kmalloc}。}它的記憶體在 kernel 的線性映射（linear map：kernel 頁表裡把整個 DDR 一對一對應進來的那一段，cacheable）裡，要給 PL 就得自己做 streaming DMA；一次最多 4 MB；也不給 \code{dma_addr_t}，而有 SMMU 時實體位址根本不是 PL 該用的位址。

\textbf{兩種「不經過 cache」的記憶體。}arm64 頁表的每一頁都帶屬性。這裡牽涉兩種，CPU 存取都不經過 cache，差別在寫入：

\begin{center}
\small
\begin{tabularx}{\textwidth}{l X X}
\toprule
 & Device-nGnRnE & Normal non-cacheable \\
\midrule
誰會得到 & \code{/dev/mem} 映射 no-map 區、\code{pgprot_noncached()} & \code{dma_mmap_coherent()}、\code{pgprot_dmacoherent()} \\
寫入 & 每個 store 都是一筆獨立的匯流排交易，不合併（nG）；回應必須由終點發出（nE） & 連續的小 store 可以在 write buffer 裡合併成一筆大的寫入（write gathering） \\
讀取 & 每次都到 DDR & 每次都到 DDR \\
對齊 & 不能非對齊存取 & 可以 \\
\bottomrule
\end{tabularx}
\end{center}

所以寫入差 3 到 5 倍、讀取一樣，和量到的一致。讀取 81 MB/s 表示每 8 bytes 約 100 ns，大約是一趟 DDR 讀取的延遲，兩種屬性都要付。只有 \code{write-in} 反映到 fps，因為寫輸入在 producer 執行緒的關鍵路徑上（寫完才能啟動 PL）；讀 heatmap 在另一個執行緒，和 PL 重疊。

\textbf{實驗的限制。}\code{uio_68} 只換位置、fps 不變，排除了「buffer 在 DDR 哪裡」；但 \code{uio_68}→\code{drv_dev} 一次換了兩件事，要分開還需要一組「driver 版、buffer 在保留區」，這次沒有做，所以只能說「輪詢位置與 buffer 來源合起來看不出影響」。

\subsection{面試會問}

\begin{qa}{為什麼用 dma\_alloc\_coherent？HP port 不是 coherent，cache 怎麼處理？driver 為什麼比較快？}
\code{dma_alloc_coherent} 一次給我 CPU 位址和裝置用的 DMA 位址。不用 \code{kmalloc}，因為它是 cacheable、最大 4 MB、也不給 DMA 位址；不用 \code{/dev/mem}，因為它能讀寫整個實體記憶體。HP port 不查 CPU 的 cache，節點沒有 \code{dma-coherent}，kernel 就把 buffer 映射成 non-cacheable。沒有 SMMU 時它從 CMA 拿連續的頁，有 SMMU 時從一般記憶體配、由 SMMU 排成連續的 IOVA。driver 版本比 UIO 版本快約 1\%（同一組開機檔上 1.1\%），我做四組 A/B 實驗，把記憶體屬性單獨拆出來，只差這一項的兩組差 1.2\%：原因是 CPU 端屬性從 Device 變成 Normal non-cacheable，寫入可以合併，CPU 寫 buffer 快 3 到 5 倍，寫輸入從 \splitDevIn{} 降到 \splitNcIn{} ms。
\end{qa}

\subsection{關鍵函式速查}

\begin{center}
\small
\begin{tabularx}{\textwidth}{l l X}
\toprule
名稱 & 位置 & 做什麼 \\
\midrule
probe 第 3 步 & \code{twn_fuse.c}:\lineof{drv_probe_dma} & 40-bit mask、配 buffer、檢查對齊 \\
\code{arena_mb} & \code{twn_fuse.c}:\lineof{drv_param_arena} & buffer 大小（MB），預設 8 \\
buffer 配置 & \code{mbv2_fuse_ps.c}:\lineof{app_arena_map} & app 在 8 MB 裡的區段 \\
\code{run()} & \code{test/ab_measure.sh}:\lineof{ab_run} & A/B 實驗：切 driver、跑 pipeline 三次 \\
第 3 部分 & \code{test/twn_bench.c}:\lineof{bench_copy} & CPU 寫、讀 buffer 的頻寬 \\
\code{dmam_alloc_coherent} & kernel & coherent buffer，driver 離開時自動歸還 \\
\bottomrule
\end{tabularx}
\end{center}
